\documentclass{beamer}

\setbeamertemplate{navigation symbols}{}

\usepackage{ctex}
\usepackage{lmodern}
%\usepackage[T1]{fontenc}
\usepackage{amsmath}

\makeatletter
% 文字上方的可伸缩双线右箭头
\newcommand{\overRightarrow}[1]{%
  \mathpalette{\overarrow@\Rightarrowfill@}{#1}%
}
% 文字上方的可伸缩双线左箭头（附赠）
\newcommand{\overLeftarrow}[1]{%
  \mathpalette{\overarrow@\Leftarrowfill@}{#1}%
}
% 文字上方的可伸缩双线双向箭头（附赠）
\newcommand{\overLeftrightarrow}[1]{%
  \mathpalette{\overarrow@\Leftrightarrowfill@}{#1}%
}
\makeatother

\usepackage{tikz}
\usepackage{tikz-cd}

\usepackage{hyperref}

%\usepackage{unicode-math}
%\setmathfont{STIX Two Math}
%\setmathfont{TeX Gyre Termes Math}
%\setmathfont{Libertinus Math}
%\setmathfont{Fira Math}

%%%%%%%% FONT TIMESAVES %%%%%%%
\newcommand{\bT}{\mathbb{T}}
\newcommand{\bN}{\mathbb{N}}
\newcommand{\bI}{\mathbb{I}}

\newcommand{\cC}{\mathcal{C}}

%%%%%%%% CATEGORY THEORY %%%%%%%
\newcommand{\id}{\mathrm{id}}
\DeclareMathOperator{\Hom}{Hom}

\newcommand{\cat}[1]{\mathrm{#1}}
\newcommand{\Set}{\cat{Set}}

\newcommand{\Id}{\text{Id}}
\newcommand{\refl}{\text{refl}}

\DeclareFontFamily{U}{dmjhira}{}
\DeclareFontShape{U}{dmjhira}{m}{n}{ <-> dmjhira }{}

\DeclareRobustCommand{\yo}{\text{\usefont{U}{dmjhira}{m}{n}\symbol{"48}}}

\DeclareMathOperator*{\colim}{colim}

\newcommand{\inl}{\operatorname{inl}}
\newcommand{\inr}{\operatorname{inr}}

%%%%%%%%% GEOMTERY %%%%%%%%%%%%%%
\newcommand{\Spec}{\operatorname{Spec}}
\newcommand{\Zar}{\mathrm{Zar}}

%%%%%%%%% TYPE THEORY %%%%%%%%%%%
\newcommand{\OO}{\bigcirc}
\newcommand{\UU}{\mathcal{U}}
\newcommand{\isEquiv}{\bibopenparen{isEquiv}}
\newcommand{\hlevel}{\mathrm{hlevel}}
\newcommand{\isContr}{\operatorname{isContr}}
\newcommand{\isProp}{\operatorname{isProp}}
\newcommand{\Prop}{\mathrm{Prop}}
%\newcommand{\refl}{\mathrm{refl}}
\newcommand{\isModal}{\operatorname{isModal}}
\newcommand{\ap}{\operatorname{ap}}
\newcommand{\fib}{\operatorname{fib}}
\newcommand{\im}{\operatorname{im}}
\newcommand{\tot}{\mathrm{tot}}
\newcommand{\pt}{\mathrm{pt}}

\title{Higher Observational Type Theory}

\date{}

\begin{document}

\begin{frame}
    % Print the title page as the first slide
    \titlepage
\end{frame}

\begin{frame}{相等性的第一原则}
对于任何类型 $A$ 和任何 $x, y : A$，存在一个恒等类型 $\Id_A(x, y)$。也就是说，$\Id_A : A \times A \to \UU$。我们为每个 $A$ 分别定义 $\Id_A$。
\end{frame}

\begin{frame}{相等性的第二原则}
对于任何类型 $A$ 和任何元素 $x : A$，存在一个自反性元素 $\refl_x : \Id_A(x, x)$。我们为每个 $x$ 分别定义 $\refl_x$。
\end{frame}

\begin{frame}{相等性的第三原则}
所有构造都尊重相等性：如果它们的输入全部被替换为相等的东西，输出也将相等。我们为每个构造分别证明/定义这一点。
\end{frame}

\begin{frame}{相等性的第四原则}
任何 $E : \Id_{\UU}(A, B)$ 产生：
\begin{itemize}
\item 一个函数 $E : A \times B \to \UU$。
\item 对于任何 $a : A$，一个元素 $\overrightarrow{E}(a) : B$。
\item 对于任何 $a : A$，一个元素 $\overRightarrow{E}(a) : E(a, \overrightarrow{E}(a))$。
\item 对于任何 $b : B$，一个元素 $\overleftarrow{E}(b) : A$。
\item 对于任何 $b : B$，一个元素 $\overLeftarrow{E}(b) : E(\overleftarrow{E}(b), b)$。
\end{itemize}
此外，对于 $A : \UU$，我们有 $\refl_A(a_0, a_1) = \Id_A(a_0, a_1)$。
\end{frame}

\begin{frame}[fragile]{相等性的第五原则}
每个方块都有一个关联的对称/转置方块：
\[
\begin{tikzcd}[row sep=large, column sep=large]
a_{10} \arrow[r, "a_{12}"] & a_{11}  \\
a_{00} \arrow[u, "a_{20}"] \arrow[r, "a_{02}"'] \arrow[ur, phantom, "a_{22}"] & a_{01} \arrow[u, "a_{21}"']
\end{tikzcd}
\quad\rightsquigarrow\quad
\begin{tikzcd}[row sep=large, column sep=large]
a_{01} \arrow[r, "a_{21}"] & a_{11}  \\
a_{00} \arrow[u, "a_{02}"] \arrow[r, "a_{20}"'] \arrow[ur, phantom, "\text{sym}(a_{22})"] & a_{10} \arrow[u, "a_{12}"']
\end{tikzcd}
\]
我们为方块的每个构造分别定义这些。
\end{frame}

\begin{frame}[fragile]{第五原则的分析}
\begin{enumerate}
  \item $\Id_{A \to B}(f_0, f_1) = \prod_{a_0, a_1 : A} \left( \Id_A(a_0, a_1) \to \Id_B(f_0(a_0), f_1(a_1)) \right).$
  \item $f : A \to B$
  \item $\refl_f : \Id_{A\to B}(f, f)$
  \item $\refl_f : \prod_{a_0, a_1 : A} \left( \Id_A(a_0, a_1) \to \Id_B(f(a_0), f(a_1)) \right).$
  \item $\Id_A : A \times A \to \UU$
  \item $\refl_{\Id_A} : \Id_{A\times A\to\UU}(\Id_A,\Id_A)$
  \item $\refl_{\Id_A} : \prod_{(x_0,y_0),(x_1,y_1):A\times A}(\Id_{A\times A}((x_0,y_0),(x_1,y_1)) \to \Id_{\UU}(\Id_A(x_0,y_0),\Id_A(x_1,y_1)))$
  \item $\Id_{A\times A}((x_0,y_0),(x_1,y_1)) = \Id_A(x_0,x_1)\times\Id_A(y_0,y_1)$
  \item 把 $\refl_{\Id_A}$ 应用到 $(a_{00},a_{10})$和$(a_{01},a_{11})$上
\end{enumerate}
\end{frame}

\begin{frame}[fragile]
  \begin{enumerate}
    \setcounter{enumi}{9}
    \item $\refl_{\Id_A}(a_{00},a_{10})(a_{01},a_{11}) : (\Id_A(a_{00},a_{01})\times\Id_A(a_{10},a_{11}))\to\Id_{\UU}(\Id_A(a_{00},a_{10}),\Id_A(a_{01},a_{11}))$
    \item 取 $a_{02}:\Id_A(a_{00},a_{01})$ 和 $a_{12} : \Id_A(a_{10},a_{11})$
    \item 把 $\refl_{\Id_A}(a_{00},a_{10})(a_{01},a_{11})(a_{02},a_{12}) : \Id_{\UU}(\Id_A(a_{00},a_{10}),\Id_A(a_{01},a_{11}))$
    \item 把 $\refl_{\Id_A}(a_{00},a_{10})(a_{01},a_{11})(a_{02},a_{12})$ 简写为 $\refl_{\Id_A}(a_{02},a_{12})$
    \item 根据第四原则, 存在 $\refl_{\Id_A}(a_{02},a_{12}) : \Id_A(a_{00},a_{10})\times\Id_A(a_{01},a_{11}) \to \UU$
    \item 取 $a_{20} : \Id_A(a_{00},a_{10})$ 和 $a_{21} : \Id_A(a_{01},a_{11})$
    \item $\refl_{\Id_A}(a_{02},a_{12})(a_{20},a_{21}) : \UU$
    \item \[
    \begin{tikzcd}[row sep=large, column sep=large]
a_{10} \arrow[r, "a_{12}"] & a_{11}  \\
a_{00} \arrow[u, "a_{20}"] \arrow[r, "a_{02}"'] \arrow[ur, phantom, "a_{22}"] & a_{01} \arrow[u, "a_{21}"']
\end{tikzcd}
    \]
  \end{enumerate}
\end{frame}

\begin{frame}{依赖函数的相等}
对 $f,g : \prod_{x:A}B(x)$, 定义
\[
\Id_{\prod_{x:A}B(x)}(f,g) = \prod_{x_0,x_1:A}\prod_{p:\Id(x_0,x_1)}\refl_B(p)(f(x_0),g(x_1))
\]
下面分析这个定义
\begin{enumerate}
  \item $B : A \to \UU$
  \item $\refl_B : \Id_{A\to\UU}(B,B)$
  \item $\refl_B : \prod_{x_0,x_1:A}(\Id_A(x_0,x_1)\to \Id_{\UU}(B(x_0),B(x_1)))$
  \item $\refl_B(p) := \refl_B(x_0,x_1,p) : \Id_{\UU}(B(x_0),B(x_1))$
  \item 根据第四原则, 存在 $\refl_B(p) : B(x_0) \times B(x_1) \to \UU$
  \item $f(x_0) : B(x_0)$, $g(x_1) : B(x_1)$
  \item $\refl_B(p)(f(x_0),g(x_1)) : \UU$
\end{enumerate}
\end{frame}

\begin{frame}[fragile]{$x\mapsto\refl_x$}
  若 $x : A$, 那么 $\refl_x : \Id_A(x,x)$, 所以
  \[
    x \mapsto \refl_x : \prod_{x:A}\Id_A(x,x)
  \]
  在上面依赖函数的定义中, 令 $f := x\mapsto \refl_x$, $B(x) := \Id_A(x,x)$
  \begin{enumerate}
    \item $\refl_{x\mapsto\refl_x} : \Id_{\prod_{x:A}\Id_A(x,x)}(x\mapsto\refl_x,x\mapsto\refl_x)$
    \item $\refl_f : \Id_{\prod_{x:A}B(x)}(f,f)$
    \item $\Id_{\prod_{x:A}B(x)}(f,f) = \prod_{x_0,x_1:A}\prod_{p:\Id(x_0,x_1)}\refl_B(p)(f(x_0),f(x_1)) = \prod_{x_0,x_1:A}\prod_{p:\Id(x_0,x_1)}\refl_{\Id_A}(p,p)(\refl_{x_0},\refl_{x_1})$
    \item $\refl_{x\mapsto\refl_x}(a_0,a_1,a_2) : \refl_{\Id_A}(a_2,a_2)(\refl_{a_0},\refl_{a_1})$
    \item 把 $\refl_{x\mapsto\refl_x}(a_0,a_1,a_2)$ 简写为 $\refl_{x\mapsto\refl_x}(a_2)$, 完全符合 square 的定义
    \item \[
    \begin{tikzcd}[row sep=large, column sep=huge]
a_{0} \arrow[r, "a_{2}"] & a_{1}  \\
a_{0} \arrow[u, "{\refl_{a_0}}"] \arrow[r, "a_{2}"'] \arrow[ur, phantom, "{\refl_{x\mapsto\refl_x}(a_2)}"] & a_{1} \arrow[u, "{\refl_{a_1}}"']
\end{tikzcd}
    \]
  \end{enumerate}

\end{frame}


% 参考文献帧，长列表自动分页
\begin{frame}[allowframebreaks]{参考文献}
\begin{thebibliography}{99}

\bibitem{txst-seminar}
\url{https://home.sandiego.edu/~shulman/papers/txst-seminar.df}

\end{thebibliography}
\end{frame}

\end{document}